\paragraph{精确几何与轨迹实验。}
二维实验独立扫描梯度比$\rho$和曲率比$\kappa$，在固定最大Hessian特征值1及共同更新RMS下比较各方向，
并沿每个方向优化标量步长。矩阵实验使用$16\times16$变量、10个配对种子，控制Hessian条件数、谱方向关系与梯度噪声。
同动量置换实验使用定理\ref{thm:v2reversal}的两个算子，逐项验证收益比和排序反转。
多步轨迹使用相同问题、初始化与噪声序列，分别筛选学习率，运行500步的余弦衰减。
梯度奇异值的对数跨度取$0,0.5,1,2$，正跨度加入固定种子的轻微扰动。
同一种子在各曲率设置中使用相同的梯度、噪声方向和随机正交基。

\paragraph{语料与模型。}
语言实验采用WikiText103\cite{merity2016wikitext}与TinyStories\cite{eldan2023tinystories}。
每个语料取固定修订版本的5000万GPT-2 BPE训练token，记录末加入结束符。
WikiText103沿用官方验证与测试划分。TinyStories将官方验证故事按内容哈希分为验证与测试两组。
模型为预归一化因果Transformer，使用GELU、学习式位置嵌入及共享输入输出权重，
上下文长度256，填充词表大小50304。
两个配置分别为4层、宽384、6头的26,499,840参数模型，和8层、宽512、8头的51,069,952参数模型。
其中隐藏矩阵参数分别为7,077,888和25,165,824。

\paragraph{训练与学习率选择。}
每次更新处理8192个token，由微批4、梯度累积8组成。前向与NS矩阵乘法采用bfloat16，参数及优化器状态采用float32。
全局梯度裁剪阈值为1，动量系数0.95，AdamW的$\beta=(0.9,0.95)$，稳定项$10^{-8}$。
二维参数权重衰减为0.1，一维参数衰减为零。辅助学习率为0.001。
学习率前10\%步数线性预热，再余弦衰减到峰值的10\%。

每个优化器族从$\{0.0003,0.001,0.003\}$开始，在种子101上运行400步。
边界最优时将该边界按3倍向外扩展，随后在种子102上确认前两名，以两种子的平均验证损失选择配置。
主实验的NS1、NS3、NS5、Probe-NS、Error-NS与重放使用NS5选定的矩阵学习率。
较长训练对照使用NS3与NS5各自在筛选阶段选定的学习率。
两个语料及两个模型配置各自执行对应优化器族的筛选。

\paragraph{主对照与扩展。}
主实验比较AdamW、幅度匹配动量、NS1、NS3、NS5、Probe-NS、阶段匹配重放及Error-NS。
幅度消融在AdamW的学习率下匹配Adam方向的RMS，在NS5的学习率下使用Moonlight尺度。
规模与语料扩展分别比较AdamW、NS5及Probe-NS。
各设置运行种子0、1、2。同一种子共享初始化及训练数据序列，主对照的方法顺序按种子轮转。
主训练运行1000次更新，处理8,192,000个token。较长训练运行2500次更新，处理20,480,000个token。

\paragraph{局部干预与评价。}
在第250、500和1000次更新提交前，保存基础模型、实际GPU候选及动量。
离线测量加载这些完整候选，以float32数学注意力后端计算Hessian--向量积。
批次$B$含1024个token，独立批次$C$含4096个token。每个检查点重复两组采样，
分别在$0.5,1,2$倍实际步长上比较预测与响应。
倍率作用于完整步长$S_k$，包括辅助参数和权重衰减分量。
精确极因子和经典三次NS作为参考。
记录的投影Hessian、梯度对齐、谱误差和候选损失组成可重放的局部几何数据。

训练每250步验证一次。WikiText103的验证覆盖完整分割，TinyStories使用固定262144个目标的验证前缀。
验证选择最佳检查点，再对完整测试分割计分。损失单位为nats/token，分块计分按目标数量加权。
统计以独立训练种子为单位，局部干预先在各训练运行内聚合，再比较预测器和方法的配对差异。

\paragraph{计时。}
实验在一张RTX 3070 Ti Laptop GPU（8\,GB显存）上运行。两个模型规格分别从冻结权重出发，
对各方法使用同一组训练输入，作4次预热后同步测量12步。
分别记录前向与反向、梯度裁剪，以及优化器计算和提交步长的时间。
策略选择步的计时包含生成候选和计算选择分数。完整训练的墙钟另计验证和锚点写盘。
